\section{Current and final regions of the repair family}
\label{sec:area_law_repair_generations}

The uniform repair induction controls every generation of the recursive
family of Definition~\ref{def:area_law_repair_stage}. This section follows the
regions from one generation to the next: which current regions survive, the
final region of each retained identifier, the terminal interiors left behind,
and the number of sites and identifiers produced in finitely many generations.

% Source: 10-geometry.tex, retained identifiers, lines 398--409;
% comparisons with later actual regions, lines 584--591.
\begin{definition}[Inclusions between generations]
    \label{def:area_law_repair_inclusions}
    \lean{TNLean.PEPS.AreaLaw.Geometry.repairNextOldInclusion,
        TNLean.PEPS.AreaLaw.Geometry.repairStageOldInclusion,
        TNLean.PEPS.AreaLaw.Geometry.repairRegionLift_eq_oldInclusion}
    \leanok
    \uses{def:area_law_repair_stage}
    With $\iota_{j,j+1}:I_j\to I_{j+1}$ the retained-identifier inclusion of
    Definition~\ref{def:area_law_repair_stage}, define $\iota_{j,j+\ell}:I_j\to I_{j+\ell}$ recursively by
    \begin{align}
        \iota_{j,j} & =\operatorname{id}_{I_j},
                    & \iota_{j,j+\ell+1}
                    & =\iota_{j+\ell,j+\ell+1}\circ\iota_{j,j+\ell}.
        \label{eq:area_law_repair_inclusions}
    \end{align}
\end{definition}

% Source: 10-geometry.tex, replacement instructions, lines 398--409;
% comparisons with later actual regions, lines 584--591.
% Revision: adc7f1241b42e322a6451854ab7e4b4c146bf78a.
\begin{lemma}[Decrease of retained current regions]
    \label{thm:area_law_repair_current_inclusion_subset}
    \lean{TNLean.PEPS.AreaLaw.Geometry.repairStage_current_oldInclusion_subset}
    \leanok
    \uses{def:area_law_repair_stage,def:area_law_repair_inclusions}
    For all $j,n\ge0$ and $i\in I_j$, $R_{\iota_{j,j+n}(i)}^{(j+n)}\subseteq R_i^{(j)}$.
\end{lemma}

\begin{proof}\leanok
    \uses{def:area_law_repair_successor_regions,def:area_law_repair_inclusions}
    Induct on $n$ using $R_{\iota_{j,j+1}(i)}^{(j+1)}=R_i^{(j)}\setminus Q_{K_j}\subseteq R_i^{(j)}$
    from (\ref{eq:area_law_repair_successor}) and (\ref{eq:area_law_repair_inclusions}).
\end{proof}

% Source: 10-geometry.tex, replacement instructions, lines 398--409;
% current regions, lines 558--561.
% Revision: adc7f1241b42e322a6451854ab7e4b4c146bf78a.
\begin{lemma}[Current regions under successive inclusions]
    \label{thm:area_law_repair_current_inclusion_add}
    \lean{TNLean.PEPS.AreaLaw.Geometry.repairStage_current_oldInclusion_add}
    \leanok
    \uses{def:area_law_repair_stage,def:area_law_repair_inclusions}
    For all $j,n,m\ge0$ and $i\in I_j$,
    \begin{align}
        R_{\iota_{j,j+(n+m)}(i)}^{(j+(n+m))}
         & =R_{\iota_{j+n,(j+n)+m}(\iota_{j,j+n}(i))}^{((j+n)+m)}.\notag
    \end{align}
\end{lemma}

\begin{proof}\leanok
    \uses{def:area_law_repair_successor_regions,def:area_law_repair_inclusions}
    Induct on $m$: $m=0$ is immediate, and at $m+1$ both sides subtract the same removal union
    $Q_{K_{j+n+m}}$ by (\ref{eq:area_law_repair_successor}) and (\ref{eq:area_law_repair_inclusions}).
\end{proof}

% Source: 10-geometry.tex, replacement instructions, lines 398--409;
% termination, lines 538--541; current regions, lines 558--561.
% Revision: adc7f1241b42e322a6451854ab7e4b4c146bf78a.
\begin{lemma}[Removal by an empty mark family]
    \label{thm:area_law_repair_removal_empty}
    \lean{TNLean.PEPS.AreaLaw.Geometry.repairRemovalSet_eq_empty}
    \leanok
    \uses{def:area_law_simultaneous_remnants}
    If $K=\varnothing$, then $\bigcup_{p\in K}Q(v_p,S_p)=\varnothing$.
\end{lemma}

\begin{proof}\leanok
    Immediate since the index set $K$ is empty.
\end{proof}

\begin{lemma}[Current regions after an empty generation]
    \label{thm:area_law_repair_current_empty_step}
    \lean{TNLean.PEPS.AreaLaw.Geometry.repairStage_current_oldInclusion_one_of_isEmpty}
    \leanok
    \uses{def:area_law_repair_stage,def:area_law_repair_inclusions}
    If $K_j=\varnothing$, then $R_{\iota_{j,j+1}(i)}^{(j+1)}=R_i^{(j)}$ for every $i\in I_j$.
\end{lemma}

\begin{proof}\leanok
    \uses{def:area_law_simultaneous_remnants,thm:area_law_repair_removal_empty}
    By (\ref{eq:area_law_repair_successor}) and Lemma~\ref{thm:area_law_repair_removal_empty},
    $R_{\iota_{j,j+1}(i)}^{(j+1)}=R_i^{(j)}\setminus\varnothing=R_i^{(j)}$.
\end{proof}

\begin{lemma}[Stability after an empty generation]
    \label{thm:area_law_repair_current_empty_later}
    \lean{TNLean.PEPS.AreaLaw.Geometry.repairStage_current_oldInclusion_of_isEmpty}
    \leanok
    \uses{def:area_law_repair_stage,def:area_law_repair_inclusions}
    If $K_j=\varnothing$, then $R_{\iota_{j,j+n}(i)}^{(j+n)}=R_i^{(j)}$ for all $n\ge0$ and $i\in I_j$.
\end{lemma}

\begin{proof}\leanok
    \uses{def:area_law_repair_ancestors,thm:area_law_repair_current_empty_step}
    Since $K_j=\varnothing$ implies $K_{j+n}=\varnothing$ for all $n\ge0$, induction on $n$ via
    Lemma~\ref{thm:area_law_repair_current_empty_step} and (\ref{eq:area_law_repair_inclusions})
    yields the identity.
\end{proof}

% Source: 10-geometry.tex, replacement instructions, lines 398--409;
% current and final regions, lines 558--561 and 584--591.
% Revision: adc7f1241b42e322a6451854ab7e4b4c146bf78a.
\begin{lemma}[Final regions at an empty generation]
    \label{thm:area_law_repair_final_current_empty}
    \lean{TNLean.PEPS.AreaLaw.Geometry.repairFinalRegion_eq_current_of_isEmpty}
    \leanok
    \uses{def:area_law_repair_stage,def:area_law_repair_inclusions}
    For $i\in I_j$, let $R_i^{(\infty,j)}=\bigcap_{n\ge0}R_{\iota_{j,j+n}(i)}^{(j+n)}$.
    If $K_j=\varnothing$, then $R_i^{(\infty,j)}=R_i^{(j)}$ for every $i\in I_j$.
\end{lemma}

\begin{proof}\leanok
    \uses{def:area_law_repair_inclusions,thm:area_law_repair_current_empty_later}
    By Lemma~\ref{thm:area_law_repair_current_empty_later}, every term in
    $\bigcap_{n\ge0}R_{\iota_{j,j+n}(i)}^{(j+n)}$ equals $R_i^{(j)}$.
\end{proof}

% Source: OpenAI2026AreaLaw, Section 11, 10-geometry.tex,
% replacement instructions, lines 398--409; final actual regions,
% lines 558--561; comparison with earlier actual regions, lines 584--591.
% Revision: adc7f1241b42e322a6451854ab7e4b4c146bf78a.
% Auxiliary consequence of the retained-region construction.
\begin{lemma}[Final regions of retained identifiers]
    \label{thm:area_law_repair_final_retention}
    \lean{TNLean.PEPS.AreaLaw.Geometry.repairFinalRegion_eq_final_oldInclusion}
    \leanok
    \uses{def:area_law_repair_stage,def:area_law_repair_inclusions}
    For every $j,n\ge0$ and $i\in I_j$, $R_i^{(\infty,j)}=R_{\iota_{j,j+n}(i)}^{(\infty,j+n)}$.
\end{lemma}

\begin{proof}\leanok
    \uses{def:area_law_repair_inclusions,
        thm:area_law_repair_current_inclusion_subset,
        thm:area_law_repair_current_inclusion_add}
    Since $A_m=R_{\iota_{j,j+m}(i)}^{(j+m)}$ decreases in $m$
    (Lemma~\ref{thm:area_law_repair_current_inclusion_subset}), Lemma~\ref{thm:area_law_repair_current_inclusion_add}
    gives $R_i^{(\infty,j)}=\bigcap_{m\ge0}A_m=\bigcap_{m\ge0}A_{n+m}=R_{\iota_{j,j+n}(i)}^{(\infty,j+n)}$.
\end{proof}

% Source: OpenAI2026AreaLaw, Section 11, 10-geometry.tex,
% replacement instructions, lines 398--409; final actual regions,
% lines 558--561; comparison with earlier actual regions, lines 584--591.
% Revision: adc7f1241b42e322a6451854ab7e4b4c146bf78a.
% Auxiliary consequence of the existing retained-region construction.
\begin{lemma}[Earlier final regions at an empty generation]
    \label{thm:area_law_repair_final_current_later_empty}
    \lean{TNLean.PEPS.AreaLaw.Geometry.repairFinalRegion_eq_current_oldInclusion_of_isEmpty}
    \leanok
    \uses{def:area_law_repair_stage,def:area_law_repair_inclusions}
    If $K_{j+n}=\varnothing$ ($j,n\ge0$), then
    $R_i^{(\infty,j)}=R_{\iota_{j,j+n}(i)}^{(j+n)}$ for every $i\in I_j$.
\end{lemma}

\begin{proof}\leanok
    \uses{thm:area_law_repair_final_retention,
        thm:area_law_repair_final_current_empty}
    Combine Lemmas~\ref{thm:area_law_repair_final_retention} and
    \ref{thm:area_law_repair_final_current_empty} at generation $j+n$.
\end{proof}

% Source: 10-geometry.tex, ordering before geometry:birth-separation,
% lines 565--571. The dummy is excluded from the physical order.
\begin{definition}[Birth rank]
    \label{def:area_law_repair_birth_rank}
    \lean{TNLean.PEPS.AreaLaw.Geometry.initialRegionBirthRank,
        TNLean.PEPS.AreaLaw.Geometry.repairStageBirthRank}
    \leanok
    \uses{def:area_law_repair_stage,def:area_law_repair_inclusions}
    Define $\operatorname{rk}_j:I_j\to\mathbb N$ by assigning $\operatorname{rk}_0(i)=0$ to
    initial primaries and the dummy $*_0$, $\operatorname{rk}_0(i)=1$ to initial belts, and for
    $i\in I_j$ and fresh $q\in I_{j+1}$,
    \begin{align}
        \operatorname{rk}_{j+1}(\iota_{j,j+1}(i))
         & =\operatorname{rk}_j(i),
         & \operatorname{rk}_{j+1}(q) & =j+2.
        \label{eq:area_law_repair_birth_rank}
    \end{align}
    The dummy is excluded from the physical order. Birth rank prescribes
    the order of the birth classes; ties within each class may be broken
    arbitrarily, as prescribed before
    \cite[Lemma~11.5]{OpenAI2026AreaLaw}.
\end{definition}

% Source: 10-geometry.tex, terminal interiors, lines 398--409;
% open actual regions and closed birth regions, lines 558--561.
% This set is not identified with the full exceptional remainder, which
% also involves the current dummy region and the treatment of interfaces.
\begin{definition}[Accumulated terminal interiors]
    \label{def:area_law_repair_terminal_interiors}
    \lean{TNLean.PEPS.AreaLaw.Geometry.repairTerminalInteriors}
    \leanok
    \uses{def:area_law_repair_stage}
    With $E_{\mathrm{term}}(\mathcal D_j)$ the terminal-square interiors at generation $j$
    (Definition~\ref{def:area_law_repair_terminal_children}), define
    \begin{align}
        \mathcal E_{\mathrm{term}}
         & =\bigcup_{j\ge0}E_{\mathrm{term}}(\mathcal D_j).
        \label{eq:area_law_repair_terminal_interiors}
    \end{align}
\end{definition}

% Source: 10-geometry.tex, geometry:bad-sites, lines 694--707.
% Revision: adc7f1241b42e322a6451854ab7e4b4c146bf78a.
\begin{lemma}[Sites in one terminal square]
    \label{thm:area_law_repair_terminal_square_sites}
    \lean{TNLean.PEPS.AreaLaw.Geometry.card_closedSquareSample_repairHalfSide_le}
    \leanok
    \uses{def:area_law_repair_square,def:peps_closed_square_sample}
    For finite $\Lambda\subset\mathbb Z^2$, $v\in\mathbb R^2$, $S>0$, and $S_*\in\mathbb R$ with
    $d(S)<S_*$, the closed square sample $K(v,r)$ (Definition~\ref{def:peps_closed_square_sample})
    satisfies $|K(v,d(S))|\le(2S_*+1)^2$.
\end{lemma}

\begin{proof}
    \leanok
    \uses{thm:peps_closed_square_card}
    Since $0<d(S)<S_*$, $K(v,d(S))\subseteq K(v,S_*)$, and Lemma~\ref{thm:peps_closed_square_card}
    at radius $S_*$ gives $|K(v,S_*)|\le(2S_*+1)^2$.
\end{proof}

% Source: 10-geometry.tex, geometry:bad-sites, lines 694--707.
% Revision: adc7f1241b42e322a6451854ab7e4b4c146bf78a.
\begin{lemma}[Terminal sites in finitely many generations]
    \label{thm:area_law_repair_terminal_prefix_sites}
    \lean{TNLean.PEPS.AreaLaw.Geometry.card_repairTerminalInteriorSites_prefix_le}
    \leanok
    \uses{def:area_law_repair_stage,def:area_law_repair_terminal_children}
    For any $o\in\mathbb R^2$, $S_*\in\mathbb R$, $C_0\ge2$, $k_0\ge50{,}000{,}000$, finite
    $\Lambda\subset\mathbb Z^2$, $A\subseteq\Lambda$, and $N\in\mathbb N$,
    \begin{align}
        \left|A\cap\bigcup_{0\le j<N}E_{\mathrm{term}}(\mathcal D_j)\right|
         & \le (2S_*+1)^2\sum_{0\le j<N}|K_j|.
        \label{eq:area_law_repair_terminal_prefix_sites}
    \end{align}
\end{lemma}

\begin{proof}
    \leanok
    \uses{thm:area_law_repair_finite_marks,
        thm:area_law_repair_terminal_square_sites}
    Each $K_j$ is finite (Lemma~\ref{thm:area_law_repair_finite_marks}) and each terminal square at
    generation $j$ contains $\le(2S_*+1)^2$ lattice sites
    (Lemma~\ref{thm:area_law_repair_terminal_square_sites}); summing over $p\in K_j$ ($d(S_p)<S_*$)
    and $0\le j<N$ yields (\ref{eq:area_law_repair_terminal_prefix_sites}).
\end{proof}
% Source: 10-geometry.tex, termination, lines 538--541;
% final regions, lines 558--561; exceptional sites, lines 694--707.
% Revision: adc7f1241b42e322a6451854ab7e4b4c146bf78a.
\begin{lemma}[Terminal interiors before an empty generation]
    \label{thm:area_law_repair_terminal_interiors_before_empty}
    \lean{TNLean.PEPS.AreaLaw.Geometry.repairTerminalInteriors_eq_iUnion_lt_of_isEmpty}
    \leanok
    \uses{def:area_law_repair_terminal_interiors}
    If $N\in\mathbb N$ satisfies $K_N=\varnothing$, then
    $\mathcal E_{\mathrm{term}}=\bigcup_{0\le j<N}E_{\mathrm{term}}(\mathcal D_j)$.
\end{lemma}

\begin{proof}\leanok
    \uses{def:area_law_repair_terminal_children,thm:area_law_repair_empty_persistence}
    By Lemma~\ref{thm:area_law_repair_empty_persistence}, $K_j=\varnothing$ and hence
    $E_{\mathrm{term}}(\mathcal D_j)=\varnothing$ for all $j\ge N$.
\end{proof}

% Source: 10-geometry.tex, replacement instructions, lines 398--417;
% geometry:repair-induction, lines 474--478 and 526--541;
% final regions, lines 558--561.
% Revision: adc7f1241b42e322a6451854ab7e4b4c146bf78a.
\begin{lemma}[Current regions avoid accumulated terminal interiors]
    \label{thm:area_law_repair_current_terminal_empty}
    \lean{TNLean.PEPS.AreaLaw.Geometry.RepairMetricHypotheses.current_disjoint_terminalInteriors_of_isEmpty}
    \leanok
    \uses{def:area_law_repair_stage,def:area_law_repair_terminal_interiors,
        lem:area_law_repair_metric_invariants}
    Under the metric hypotheses of Lemma~\ref{lem:area_law_repair_metric_invariants}, if
    $K_N=\varnothing$ ($N\in\mathbb N$), then $R_i^{(N)}\cap\mathcal E_{\mathrm{term}}=\varnothing$
    for every $i\in I_N$.
\end{lemma}

\begin{proof}
    \leanok
    \uses{thm:area_law_repair_terminal_interiors_before_empty,
        def:area_law_repair_terminal_children,def:area_law_repair_ancestors,
        thm:area_law_repair_current_terminal_disjoint}
    By Lemma~\ref{thm:area_law_repair_terminal_interiors_before_empty},
    $\mathcal E_{\mathrm{term}}=\bigcup_{j<N}E_{\mathrm{term}}(\mathcal D_j)\subseteq\mathcal E_N$,
    and $R_i^{(N)}\cap\mathcal E_N=\varnothing$ by Lemma~\ref{thm:area_law_repair_current_terminal_disjoint}.
\end{proof}

% Source: 10-geometry.tex, positive boundary, lines 102--105;
% replacement instructions, lines 398--417; lattice partition, lines 545--561.
% Revision: adc7f1241b42e322a6451854ab7e4b4c146bf78a.
\begin{lemma}[Unique final-region membership outside terminal interiors]
    \label{thm:area_law_repair_final_lattice_partition_empty}
    \lean{TNLean.PEPS.AreaLaw.Geometry.RepairMetricHypotheses.lattice_mem_terminal_or_unique_final_of_isEmpty}
    \leanok
    \uses{def:area_law_repair_stage,def:area_law_repair_inclusions,
        def:area_law_dyadic_origin,def:area_law_repair_terminal_interiors,
        lem:area_law_repair_metric_invariants}
    At $o=(\sqrt2,\sqrt3)$, under the hypotheses of Lemma~\ref{lem:area_law_repair_metric_invariants}
    with $S_*\ge1$, $b=|\partial_\Lambda A|>0$, and $K_N=\varnothing$ ($N\in\mathbb N$), every
    $z\in\mathbb Z^2$ satisfies exactly one of:
    \begin{enumerate}
        \item $z\in\mathcal E_{\mathrm{term}}$ and $z\notin R_i^{(\infty,N)}$ for all $i\in I_N$;
        \item $z\notin\mathcal E_{\mathrm{term}}$ and there exists a unique $i\in I_N$ with
              $z\in R_i^{(\infty,N)}$.
    \end{enumerate}
\end{lemma}

\begin{proof}
    \leanok
    \uses{thm:area_law_repair_current_terminal_empty,
        thm:area_law_repair_final_current_empty,
        thm:area_law_repair_current_disjoint,
        thm:area_law_repair_lattice_cover_before,
        def:area_law_repair_terminal_interiors}
    Since $R_i^{(\infty,N)}=R_i^{(N)}$ (Lemma~\ref{thm:area_law_repair_final_current_empty}),
    Lemma~\ref{thm:area_law_repair_current_terminal_empty} gives (1) when
    $z\in\mathcal E_{\mathrm{term}}$, while for $z\notin\mathcal E_{\mathrm{term}}$
    Lemmas~\ref{thm:area_law_repair_lattice_cover_before} and
    \ref{thm:area_law_repair_current_disjoint} give existence and uniqueness in (2).
\end{proof}

% Source: 10-geometry.tex, geometry:trim-radius, lines 393--399.
\begin{lemma}[Closed removal sets]
    \label{thm:area_law_repair_removal_closed}
    \lean{TNLean.PEPS.AreaLaw.Geometry.isClosed_repairRemovalSet}
    \leanok
    \uses{def:area_law_repair_square,def:area_law_simultaneous_remnants}
    For any finite $K$, $Q_K=\bigcup_{p\in K}Q(v_p,S_p)$ is closed.
\end{lemma}
\begin{proof}\leanok
    Finite union of closed squares $Q(v_p,S_p)$.
\end{proof}

% Source: 10-geometry.tex, replacement instructions, lines 393--417;
% remembered closed births and current regions, lines 558--561.
\begin{lemma}[Current regions remain inside their births]
    \label{thm:area_law_repair_current_birth_containment}
    \lean{TNLean.PEPS.AreaLaw.Geometry.repairStage_current_subset_birth}
    \leanok
    \uses{def:area_law_repair_stage,def:area_law_repair_successor_regions}
    For every $j\ge0$ and $i\in I_j$,
    \begin{align}
        R_i^{(j)} & \subseteq B_i^{(j)}.
        \label{eq:area_law_repair_current_birth_containment}
    \end{align}
\end{lemma}
\begin{proof}\leanok
    By induction on $j$: at $j=0$, $R_i^{(0)}=\operatorname{int}B_i^{(0)}\subseteq B_i^{(0)}$
    by (\ref{eq:area_law_repair_initial}); at $j+1$, (\ref{eq:area_law_repair_successor}) gives
    $R_i^{(j+1)}=R_i^{(j)}\setminus Q_{K_j}\subseteq B_i^{(j)}=B_i^{(j+1)}$ for old $i$ and
    $R_q^{(j+1)}=\operatorname{int}B_q^{(j+1)}\subseteq B_q^{(j+1)}$ for fresh $q$.
\end{proof}

% Source: 10-geometry.tex, initial regions, lines 172--177,
% 212--218 and 313--316; replacement instructions, lines 393--417;
% open actual regions, lines 558--561.
% Revision: adc7f1241b42e322a6451854ab7e4b4c146bf78a.
\begin{lemma}[Openness of current regions]
    \label{thm:area_law_repair_current_open}
    \lean{TNLean.PEPS.AreaLaw.Geometry.isOpen_repairStage_current}
    \leanok
    \uses{def:area_law_repair_stage}
    For every $j\ge0$ and $i\in I_j$, $R_i^{(j)}$ is open.
\end{lemma}

\begin{proof}\leanok
    \uses{def:area_law_initial_open_regions,
        def:area_law_repair_successor_regions,thm:area_law_repair_finite_marks,
        thm:area_law_repair_removal_closed}
    Induct on $j$: $R_i^{(0)}$ is an interior, and at $j+1$ subtracting the closed set $Q_{K_j}$
    (Lemmas~\ref{thm:area_law_repair_finite_marks} and \ref{thm:area_law_repair_removal_closed})
    from open $R_i^{(j)}$ or taking $\operatorname{int}B_q^{(j+1)}$ yields open sets.
\end{proof}

% Source: 10-geometry.tex, actual regions and dummy exclusion, lines 558--568;
% final restriction to the cut, lines 777--789.
\begin{definition}[Lattice sites in the current regions]
    \label{def:area_law_repair_region_sites}
    \lean{TNLean.PEPS.AreaLaw.Geometry.RepairStage.regionSites,
        TNLean.PEPS.AreaLaw.Geometry.RepairStage.RetainedRegionIndex}
    \leanok
    \uses{def:area_law_repair_stage}
    For a stage $\mathcal D$ with current regions $(R_i)_{i\in I}$ and dummy $*$, and a finite
    cut $A\subseteq\Lambda$, define the finite site sets and retained index set
    \begin{align}
        X_i(A) & = A\cap R_i,
               & I_A          & = \{i\in I\setminus\{*\}:X_i(A)\ne\varnothing\}.
        \label{eq:area_law_repair_region_sites}
    \end{align}
\end{definition}

% Source: 10-geometry.tex, initial dummy region, lines 172--177;
% geometry:bad-sites, lines 694--707.
% Revision: adc7f1241b42e322a6451854ab7e4b4c146bf78a.
\begin{lemma}[Counting surviving dummy sites]
    \label{thm:area_law_repair_dummy_sites_count}
    \lean{TNLean.PEPS.AreaLaw.Geometry.card_repairStage_dummySites_le}
    \leanok
    \uses{def:area_law_repair_region_sites}
    With $b=|\partial_\Lambda A|$, every generation $j\ge0$ satisfies
    \begin{align}
        |X_{*_j}(A)|
         & \le 2\bigl(2(C_0+1)2^{k_0}+1\bigr)^2 b.
        \label{eq:area_law_repair_dummy_sites_count}
    \end{align}
\end{lemma}

\begin{proof}\leanok
    \uses{thm:area_law_repair_current_birth_containment,
        thm:area_law_neighborhood_endpoint_distance,
        thm:area_law_cell_counts}
    By (\ref{eq:area_law_repair_current_birth_containment}) and
    Lemma~\ref{thm:area_law_neighborhood_endpoint_distance}, every $x\in X_{*_j}(A)\subseteq\overline{N_{k_0}}$
    lies within $\ell_\infty$-distance $r=(C_0+1)2^{k_0}$ of some $z\in Z$ ($|Z|\le2b$), and each
    such ball contains $(2r+1)^2$ lattice sites (Lemma~\ref{thm:area_law_cell_counts}).
\end{proof}

% Source: 10-geometry.tex, final lattice partition, lines 558--561;
% geometry:bad-sites, lines 694--707.
\begin{definition}[Exceptional lattice sites]
    \label{def:area_law_repair_exceptional_sites}
    \lean{TNLean.PEPS.AreaLaw.Geometry.repairStageExceptionalSites}
    \leanok
    \uses{def:area_law_repair_region_sites,
        def:area_law_repair_terminal_interiors}
    At generation $N$, define the exceptional lattice sites of $A$ by
    \begin{align}
        D_N(A)
         & = A\cap\bigl(R_{*_N}^{(N)}\cup\mathcal E_{\mathrm{term}}\bigr).
        \label{eq:area_law_repair_exceptional_sites}
    \end{align}
\end{definition}

% Source: openai/math@adc7f1241b42e322a6451854ab7e4b4c146bf78a,
% Section 11, geometry:bad-sites, lines 694--707.
\begin{lemma}[Exceptional sites at an empty generation]
    \label{thm:area_law_repair_exceptional_sites_prefix_count}
    \lean{TNLean.PEPS.AreaLaw.Geometry.card_repairStageExceptionalSites_le_prefix}
    \leanok
    \uses{def:area_law_repair_exceptional_sites,def:area_law_hamiltonian}
    If $K_N=\varnothing$, then with $C_{\mathrm{dum}}=2\bigl(2(C_0+1)2^{k_0}+1\bigr)^2$ and
    $b=|\partial_\Lambda A|$,
    \begin{align}
        |D_N(A)|
         & \le C_{\mathrm{dum}}b+(2S_*+1)^2\sum_{0\le j<N}|K_j|.
        \label{eq:area_law_repair_exceptional_sites_prefix_count}
    \end{align}
\end{lemma}

\begin{proof}
    \leanok
    \uses{thm:area_law_repair_terminal_interiors_before_empty,
        thm:area_law_repair_dummy_sites_count,
        thm:area_law_repair_terminal_prefix_sites}
    By Lemma~\ref{thm:area_law_repair_terminal_interiors_before_empty},
    $D_N(A)=X_{*_N}(A)\cup\bigl(A\cap\bigcup_{0\le j<N}E_{\mathrm{term}}(\mathcal D_j)\bigr)$,
    so (\ref{eq:area_law_repair_dummy_sites_count}) and
    (\ref{eq:area_law_repair_terminal_prefix_sites}) yield
    (\ref{eq:area_law_repair_exceptional_sites_prefix_count}).
\end{proof}

% Source: 10-geometry.tex, finiteness after restriction to the cut,
% lines 777--780.
\begin{lemma}[Finiteness of the retained identifiers]
    \label{lem:area_law_repair_region_sites_finite}
    \lean{TNLean.PEPS.AreaLaw.Geometry.RepairStage.finite_retainedRegionIndex}
    \leanok
    \uses{def:area_law_repair_region_sites}
    If $(R_i)_{i\in I}$ are pairwise disjoint, then $I_A$ is finite for every finite $A$.
\end{lemma}
\begin{proof}\leanok
    The nonempty sets $(X_i(A))_{i\in I_A}$ are pairwise disjoint subsets of the finite set $A$.
\end{proof}

% Source: 10-geometry.tex, birth order, lines 565--568;
% final restriction to the cut and finite partition, lines 777--780.
% Revision: adc7f1241b42e322a6451854ab7e4b4c146bf78a.
\begin{lemma}[Ordering the retained regions at an empty generation]
    \label{thm:area_law_repair_ordered_partition_empty}
    \lean{TNLean.PEPS.AreaLaw.Geometry.RepairMetricHypotheses.exists_orderedTwoFamilyPartition_of_isEmpty}
    \leanok
    \uses{def:area_law_two_family_partition,def:area_law_repair_region_sites,
        def:area_law_repair_exceptional_sites,def:area_law_repair_birth_rank,
        def:area_law_dyadic_origin,lem:area_law_repair_metric_invariants}
    At $o=(\sqrt2,\sqrt3)$, under the hypotheses of Lemma~\ref{lem:area_law_repair_metric_invariants}
    with $S_*\ge1$, $|\partial_\Lambda A|>0$, and $K_N=\varnothing$ ($N\in\mathbb N$), there exist
    an ordered two-family partition (Definition~\ref{def:area_law_two_family_partition}) of $A$
    with pieces $(X_i)_{0\le i<m}$, labels $(f_i)_{0\le i<m}$, residual $D=D_N(A)$, and a bijection
    $e:\{0,\ldots,m-1\}\to I_A^{(N)}$ such that
    \begin{align}
        X_i    & = A\cap R_{e(i)}^{(N)},                                         &
        f_i    & = \xi_{e(i)}^{(N)},                                             &
        i\le j & \implies \operatorname{rk}_N(e(i))\le\operatorname{rk}_N(e(j)).
        \label{eq:area_law_repair_ordered_partition_empty}
    \end{align}
\end{lemma}

\begin{proof}
    \leanok
    \uses{lem:area_law_repair_region_sites_finite,
        thm:area_law_repair_current_disjoint,
        thm:area_law_repair_final_lattice_partition_empty,
        thm:area_law_repair_final_current_empty,
        thm:area_law_repair_current_terminal_empty}
    By Lemma~\ref{thm:area_law_repair_current_disjoint} and
    Lemma~\ref{lem:area_law_repair_region_sites_finite}, $I_A^{(N)}$ is finite; choose a bijection
    $e:\{0,\ldots,m-1\}\to I_A^{(N)}$ with $\operatorname{rk}_N\circ e$ nondecreasing.
    Lemmas~\ref{thm:area_law_repair_final_current_empty},
    \ref{thm:area_law_repair_current_terminal_empty}, and
    \ref{thm:area_law_repair_final_lattice_partition_empty} then show that $(X_i)_{0\le i<m}$ and
    $D=D_N(A)$ form an ordered two-family partition of $A$ satisfying
    (\ref{eq:area_law_repair_ordered_partition_empty}).
\end{proof}

% Source: 10-geometry.tex, replacement instructions, lines 393--417;
% geometry:total-repairs, lines 682--691; restriction to A, lines 777--789.
% Revision: adc7f1241b42e322a6451854ab7e4b4c146bf78a.
\begin{lemma}[Counting retained regions in finitely many generations]
    \label{thm:area_law_repair_retained_prefix_count}
    \lean{TNLean.PEPS.AreaLaw.Geometry.finite_and_card_repairStage_retainedRegionIndex_le}
    \leanok
    \uses{def:area_law_repair_stage,def:area_law_repair_region_sites}
    For every $N\in\mathbb N$, $I_A^{(N)}$ (Definition~\ref{def:area_law_repair_region_sites}) is
    finite and satisfies
    \begin{align}
        |I_A^{(N)}| & \le |I_A^{(0)}|+8\sum_{0\le j<N}|K_j|.
        \label{eq:area_law_repair_retained_prefix_count}
    \end{align}
\end{lemma}

\begin{proof}
    \leanok
    \uses{thm:area_law_initial_open_disjoint,
        lem:area_law_repair_region_sites_finite,
        thm:area_law_repair_finite_marks,thm:area_law_cell_fan_run_count,
        def:area_law_repair_successor_regions}
    $I_A^{(0)}$ is finite by Lemma~\ref{thm:area_law_initial_open_disjoint} and
    Lemma~\ref{lem:area_law_repair_region_sites_finite}. At step $j$, old regions shrink while at
    most $8|K_j|$ fresh identifiers are born (Lemmas~\ref{thm:area_law_repair_finite_marks} and
    \ref{thm:area_law_cell_fan_run_count}), so $|I_A^{(j+1)}|\le|I_A^{(j)}|+8|K_j|$.
\end{proof}

% Source: 10-geometry.tex, birth classes, lines 565--568;
% geometry:total-repairs, lines 685--692; nonprimary errors, lines 796--798.
% Revision: adc7f1241b42e322a6451854ab7e4b4c146bf78a.
\begin{lemma}[Counting identifiers of positive birth rank]
    \label{thm:area_law_repair_positive_rank_prefix_count}
    \lean{TNLean.PEPS.AreaLaw.Geometry.exists_uniform_repairStage_positive_rank_prefix_bound}
    \leanok
    \uses{def:area_law_hamiltonian,def:area_law_repair_birth_rank}
    With $B>0$ from Lemma~\ref{thm:area_law_initial_belt_count} and $b=|\partial_\Lambda A|$,
    for every $N\in\mathbb N$, the set $I_N^+=\{i\in I_N:\operatorname{rk}_N(i)>0\}$ is finite and satisfies
    \begin{align}
        |I_N^+| & \le B(2C_0+1)^2b+8\sum_{0\le j<N}|K_j|.
        \label{eq:area_law_repair_positive_rank_prefix_count}
    \end{align}
\end{lemma}

\begin{proof}
    \leanok
    \uses{thm:area_law_initial_belt_count,
        def:area_law_repair_birth_rank,
        def:area_law_repair_successor_regions,
        thm:area_law_repair_finite_marks,thm:area_law_cell_fan_run_count}
    At $N=0$, $I_0^+=\mathcal I_{\mathrm{belt}}$ has $|I_0^+|\le B(2C_0+1)^2b$ by
    Lemma~\ref{thm:area_law_initial_belt_count}; at step $j$, at most $8|K_j|$ fresh identifiers
    (of rank $j+2>0$) are added by Definition~\ref{def:area_law_repair_successor_regions} and
    Lemma~\ref{thm:area_law_cell_fan_run_count}, so $|I_{j+1}^+|\le|I_j^+|+8|K_j|$.
\end{proof}

\begin{lemma}[A uniform count of retained positive-rank identifiers]
    \label{thm:area_law_repair_retained_positive_rank_count}
    \lean{TNLean.PEPS.AreaLaw.Geometry.exists_uniform_repairStage_retained_positive_rank_bound}
    \leanok
    \uses{def:area_law_hamiltonian,def:area_law_repair_birth_rank,
        def:area_law_repair_region_sites,
        lem:area_law_repair_metric_invariants}
    For every $S_*>1$, there exists $L>0$ such that under the hypotheses of
    Lemma~\ref{lem:area_law_repair_metric_invariants}, for every $N\in\mathbb N$, the set
    $I_{A,N}^+=\{i\in I_A^{(N)}:\operatorname{rk}_N(i)>0\}$ is finite and satisfies
    \begin{align}
        |I_{A,N}^+| & \le L(2C_0+1)^2|\partial_\Lambda A|.
        \label{eq:area_law_repair_retained_positive_rank_count}
    \end{align}
\end{lemma}

\begin{proof}
    \leanok
    \uses{thm:area_law_repair_positive_rank_prefix_count,
        thm:area_law_repair_total_mark_prefix_count}
    Since $I_{A,N}^+\subseteq I_N^+$, combining (\ref{eq:area_law_repair_positive_rank_prefix_count})
    of Lemma~\ref{thm:area_law_repair_positive_rank_prefix_count} with
    (\ref{eq:area_law_repair_total_mark_prefix_count}) of
    Lemma~\ref{thm:area_law_repair_total_mark_prefix_count} gives
    (\ref{eq:area_law_repair_retained_positive_rank_count}) with $L=B+8M$.
\end{proof}

\begin{lemma}[The nonprimary error sum at a finite generation]
    \label{thm:area_law_repair_retained_nonprimary_error_sum}
    \lean{TNLean.PEPS.AreaLaw.Geometry.exists_uniform_repairStage_retained_positive_rank_error_sum_bound}
    \leanok
    \uses{def:area_law_hamiltonian,def:area_law_repair_characteristic_length,
        def:area_law_repair_region_sites,def:area_law_repair_birth_rank,
        lem:area_law_repair_metric_invariants}
    For every $S_*>1$, with $L>0$ from Lemma~\ref{thm:area_law_repair_retained_positive_rank_count},
    under the hypotheses of Lemma~\ref{lem:area_law_repair_metric_invariants}, for any $C_B>0$,
    $N\in\mathbb N$, and $n_i=\lceil C_B\tau_N(i)\rceil$ ($i\in I_{A,+}^{(N)}=I_{A,N}^+$),
    \begin{align}
        \sum_{i\in I_{A,+}^{(N)}}(n_i^{100})^{-1}
         & \le L(2C_0+1)^2|\partial_\Lambda A|.
        \label{eq:area_law_repair_retained_nonprimary_error_sum}
    \end{align}
\end{lemma}

\begin{proof}
    \leanok
    \uses{thm:area_law_repair_characteristic_length_positive,
        thm:area_law_repair_retained_positive_rank_count}
    For every $i\in I_{A,+}^{(N)}$, $i\ne *_N$ gives $\tau_N(i)>0$
    (Lemma~\ref{thm:area_law_repair_characteristic_length_positive}), so $n_i=\lceil C_B\tau_N(i)\rceil\ge1$
    and $(n_i^{100})^{-1}\le1$; summing over $I_{A,+}^{(N)}$ and applying
    (\ref{eq:area_law_repair_retained_positive_rank_count}) gives
    (\ref{eq:area_law_repair_retained_nonprimary_error_sum}).
\end{proof}
